\documentclass[11pt]{article}
\usepackage[T1]{fontenc}
\usepackage{lmodern,amsmath,amssymb,amsthm,mathtools,booktabs,microtype}
\usepackage[margin=1in]{geometry}
\usepackage[colorlinks=true,linkcolor=blue!50!black,urlcolor=blue!50!black,citecolor=blue!50!black]{hyperref}
\usepackage{xcolor}
\setlength{\emergencystretch}{3em}
\newtheorem{theorem}{Theorem}[section]
\newtheorem{lemma}[theorem]{Lemma}
\newtheorem{corollary}[theorem]{Corollary}
\newtheorem{proposition}[theorem]{Proposition}
\theoremstyle{remark}\newtheorem{remark}[theorem]{Remark}
\newcommand{\F}{\mathbb F}
\newcommand{\Z}{\mathbb Z}
\newcommand{\C}{\mathbb C}
\newcommand{\E}{\mathbb E}
\newcommand{\supp}{\operatorname{supp}}
\newcommand{\ord}{\operatorname{ord}}
\newcommand{\prm}{\operatorname{PRM}}
\newcommand{\abs}[1]{\lvert#1\rvert}
\newcommand{\norm}[1]{\lVert#1\rVert}
\newcommand{\ind}{\mathbf 1}
\title{Sharp Euler-factor agreement for cubic and quintic Dirichlet packets\\[4pt]\large Cyclotomic divisibility, projective codes, and effective detection under Dirichlet GRH}
\author{Research note prepared for the ProveIt project}
\date{27 November 2026}
\begin{document}
\maketitle
\begin{abstract}
For two distinct equal-size multisets of characters in a rank-$r$ elementary abelian group of prime exponent $\ell$, we study the density of primes where their products of local Euler factors disagree. Write $k$ for the packet size, $d=\lfloor\log_\ell k\rfloor$ and $s=\lfloor k/\ell^d\rfloor$. If $d\le r-2$, we prove the unconditional lower bound $(\ell-1)(\ell+1-s)/\ell^{d+2}$. It is sharp whenever $s$ divides $\ell-1$ or $\ell+1$, so in particular it is sharp for every size and rank at $\ell=3,5$. When $d\ge r-1$ the exact minimum is $(\ell-1)/\ell^r$. The proof combines the augmentation lattice in $\Z[\zeta_\ell]$, a Newton-identity bound on cyclotomic divisibility, and the minimum distance of projective Reed--Muller codes. We classify the support of every equality case, describe its reduction to a rank-two problem, and give explicit extremizers on the split and nonsplit conics. The Dirichlet GRH theorem exported by the vendored Lean proof yields distinguishing primes in every interval $(x/2,x)$ once $x\gg k^2\log^2(2Q)$, with a rank-independent quantitative lower bound for the weighted number of disagreements. The geometric density classification needs no GRH. This answers the minimum-density part of an explicit odd-prime question in the merged ProveIt article; it does not classify the entire spectrum of larger densities.
\end{abstract}
\tableofcontents

\section{Source theorem, motivation, and status}
The two inputs are the readable vendored Lean development and the merged research article in ProveIt \cite{vendor,article}. The vendor README identifies upstream commit \texttt{72fead102f663d5c7a407}\allowbreak\texttt{964a13dc0cf9d977207}, dated 25 November 2026, and a separate Lean 4.35.0 package. Its \texttt{src/Statements.lean} defines, in Mathlib's vocabulary, nonvanishing of all Dirichlet $L$-functions in $\Re s>1/2$ apart from the principal-character pole, and the statement that their zeros in $0<\Re s<1$ lie on $\Re s=1/2$. In \texttt{src/Main.lean} these are exported as \nolinkurl{LFunction_ne_zero_of_half_lt_re} and \nolinkurl{LFunction_eq_zero_re}; the former is obtained from \texttt{New.grh}, and \texttt{New.zeros\_on\_line} supplies the functional-equation step. The same file exports ordinary RH and the primitive-character specialization. The readable-edition README reports independent machine checks and only the standard axioms \texttt{propext}, \texttt{Classical.choice}, and \texttt{Quot.sound}. We use those statements as input, not as an independent replay of the Lean proof. The scope is Dirichlet GRH, not GRH for arbitrary Artin, Hecke, or automorphic functions.

The merged article's README describes 131 reports and 914 result headings. Its introduction distinguishes unconditional results, GRH consequences, and reductions. In \texttt{sections/rec.tex}, the theorem labelled \texttt{rec:thm:qmatch} gives the minimum disagreement density and the initial high-agreement spectrum for quadratic packets. The neighbouring \texttt{rec:thm:qmatchfinite} gives effective prime detection under GRH, and \texttt{rec:lem:gram} is its analytic estimate. The ``Pooled zero statistics'' item of \texttt{rec:sec:open} explicitly asks for analogous results for characters of odd prime order. The present note replaces the binary Walsh valuation argument by a cyclotomic augmentation-lattice argument. A key distinction is that for odd $\ell$ the disagreement set consists of punctured lines through the origin, so the relevant code is \emph{projective}, not affine.

\paragraph{Verification and novelty.} The algebraic proofs below are written out, using standard facts on cyclotomic integers, Newton identities and, for the equality characterization, a published projective-code theorem. The minimum-distance lower bound also has a self-contained proof. The analytic section uses the usual smoothed explicit formula and the vendored Dirichlet GRH statement. The Lean development was inspected at its public interfaces and the relevant article sections were read; neither the full source proof nor the hundreds of article results were rechecked. No new Lean theorem is claimed kernel-verified here. The exact minimum law is a proposed new extension relative to the cited article; priority in all mathematical literature is not asserted. A targeted literature search found the established coding-theoretic input \cite{GL,Sorensen}, not an independent statement of this packet theorem. The algebraic result is independent of the ultimate publication or recognition status of the GRH proof.

\section{Statement of the main theorem}
Let $\ell$ be a prime (we emphasize odd $\ell$) and $X$ a subgroup of Dirichlet characters modulo $Q$ isomorphic to $\F_\ell^r$, $r\ge1$. Characters are allowed to be principal or imprimitive; their Euler factors at primes not dividing $Q$ are all that is needed. The dual group $V=\widehat X$ is identified with $\F_\ell^r$ after fixing a primitive $\ell$th root $\zeta$ and a perfect pairing $a\cdot v$.
For two multisets $A,C$ of $k$ members of $X$, define
\[
 P_A(v,T)=\prod_{a\in A}(1-\zeta^{a\cdot v}T),\qquad
 c_a=\operatorname{mult}_A(a)-\operatorname{mult}_C(a),\qquad
 H(v)=\sum_{a\in X}c_a\zeta^{a\cdot v}.
\]
The signed multiplicities satisfy $\sum c_a=0$. Cancel the common factors, and let $m=\sum_{c_a>0}c_a=\frac12\sum_a\abs{c_a}$; if $A\ne C$, then $1\le m\le k$. Put $D=\{v:P_A(v,T)\ne P_C(v,T)\}$ and $\delta=\abs D/\ell^r$.

\begin{theorem}[Sharp tiers]\label{thm:main}
Let $A\ne C$ and $k\ge1$, and write $k= s\ell^d+u$ with $d=\lfloor\log_\ell k\rfloor$, $1\le s\le\ell-1$, $0\le u<\ell^d$. If $d\le r-2$, then
\begin{equation}\label{eq:main}
 \delta\ge \Delta_{\ell,d,s}:=\frac{(\ell-1)(\ell+1-s)}{\ell^{d+2}}.
\end{equation}
If $s\mid(\ell-1)$ or $s\mid(\ell+1)$, this bound is attained by packets of size $k$ in every such $X$. For $d\ge r-1$, the minimum is exactly $(\ell-1)/\ell^r$. In particular, for $\ell=3$ and $\ell=5$ every tier is sharp, for every rank and packet size. These statements hold unconditionally for the natural density of primes at which the local factors differ.
\end{theorem}

For rank $r\ge d+2$, the cubic minimum is $2/3^{d+1}$ when $s=1$ and $4/3^{d+2}$ when $s=2$. The quintic minimum, with $5^d\le k<5^{d+1}$, is
\begin{equation}\label{eq:quintic}
\frac{4}{5^{d+2}}\times
\begin{cases}5&5^d\le k<2\cdot5^d,\\4&2\cdot5^d\le k<3\cdot5^d,\\3&3\cdot5^d\le k<4\cdot5^d,\\2&4\cdot5^d\le k<5^{d+1}.\end{cases}
\end{equation}
A cubic example is $k=6$: the largest agreement density is $23/27$ if $r\ge3$, and $7/9$ if $r=2$. For $k=1000$ and $r\ge8$, $d=6,s=1$, the largest cubic agreement density is $2185/2187$.

\section{Cyclotomic reduction and the valuation bound}
\begin{lemma}[A trace suffices]\label{lem:trace}
$D=\supp H$, $H(0)=0$, and $D$ is a union of nonzero projective lines (each with $\ell-1$ points). Moreover $\delta$ is the natural density of the disagreeing primes.
\end{lemma}
\begin{proof}
For $v\in V$, put $b_j(v)=\sum_{a:a\cdot v=j}c_a$, $0\le j<\ell$. Then $\sum_j b_j=0$ and $H(v)=\sum_j b_j\zeta^j$. Because the minimal polynomial of $\zeta$ is $1+z+\cdots+z^{\ell-1}$, $H(v)=0$ forces all $b_j$ equal, hence all zero. This is exactly equality of the local polynomials, by uniqueness of factorization into their linear factors. Also $H(0)=0$, and $H(\lambda v)$ is the Galois conjugate $\zeta\mapsto\zeta^\lambda$ of $H(v)$ for $\lambda\in\F_\ell^\times$. Fourier inversion gives $H\not\equiv0$ when $c\ne0$. Evaluation from $(\Z/Q\Z)^\times$ to $V$ is onto: otherwise a nontrivial character in $X$ would be trivial on every unit. All fibers therefore have equal size, and the prime number theorem in arithmetic progressions gives density $\abs D/\ell^r$. Removing the finitely many primes dividing $Q$ changes no density.
\end{proof}

Use the valuation $v_\pi$ at $\pi=\zeta-1$ in $\mathcal O=\Z[\zeta]$, normalized by $v_\pi(\pi)=1$ and $v_\pi(\ell)=\ell-1$. We use the elementary cyclotomic facts $\mathcal O/(\pi)=\F_\ell$ and $(\ell)=(\pi^{\ell-1})$, which follow from $\prod_{j=1}^{\ell-1}(1-\zeta^j)=\ell$ and the fact that $(1-\zeta^j)/(1-\zeta)$ is a unit (its norm is one). Set $t=\min_{H(v)\ne0}v_\pi(H(v))$.

\begin{lemma}[The augmentation lattice]\label{lem:aug}
The map $b\mapsto\sum_{j=0}^{\ell-1}b_j\zeta^j$ from $\{b\in\Z^\ell:\sum b_j=0\}$ to $\pi\mathcal O$ is an isomorphism of abelian groups. In particular if $v_\pi(H(v))\ge (\ell-1)d+1$, every $b_j(v)$ is divisible by $\ell^d$.
\end{lemma}
\begin{proof}
The domain has basis $e_j-e_0$ for $1\le j<\ell$. Its images $\zeta^j-1$ equal $\pi(1+\zeta+\cdots+\zeta^{j-1})$. The latter factors after division by $\pi$ are a triangular integral basis of $\mathcal O$ in the power basis $1,\zeta,\ldots,\zeta^{\ell-2}$. Divisibility by $\pi^{(\ell-1)d+1}$ is exactly membership in $\ell^d\pi\mathcal O$, giving the assertion.
\end{proof}

\begin{lemma}[Newton bound]\label{lem:val}
With $d,s$ as in Theorem~\ref{thm:main},
\begin{equation}\label{eq:valuation}
 1\le t\le (\ell-1)d+s.
\end{equation}
The same bound holds for the valuation of each individual nonzero $H(v)$.
\end{lemma}
\begin{proof}
The lower bound is automatic since $H(v)\equiv\sum c_a=0\pmod\pi$. Suppose a nonzero value is divisible by $\pi^{(\ell-1)d+s+1}$. By Lemma~\ref{lem:aug}, $b_j=\ell^d B_j$. Then $B\ne0$, $\sum B_j=0$, and $M=\sum_{B_j>0}B_j\le k/\ell^d$, hence $M\le s<\ell$. If $s=\ell-1$, divisibility gives $\sum B_j\zeta^j\in\pi^\ell\mathcal O=\ell\pi\mathcal O$, so every $B_j$ is divisible by $\ell$, impossible. If $s<\ell-1$, expand
\[
 \sum_j B_j(1+\pi)^j=\sum_{h=1}^{\ell-1}\pi^h \sum_j B_j\binom jh.
\]
Were some coefficient at $h\le s$ nonzero modulo $\ell$, take the smallest such $h$. Its term has valuation $h$; all earlier terms have valuation at least $\ell$ and all later terms have valuation greater than $h$. This contradicts the assumed divisibility by $\pi^{s+1}$. Thus $\sum_j B_j\binom jh=0$ in $\F_\ell$ for $1\le h\le s$, and triangular change of basis (factorials are units) implies $\sum_j B_j j^h=0$ for $1\le h\le s$. The positive and negative parts of $B$ are two disjoint multisets of $M\le s$ elements of $\F_\ell$ with identical first $M$ power sums. Newton's identities, dividing only by $1,\ldots,M$, imply identical elementary symmetric polynomials and identical monic root polynomials. Hence the multisets coincide, a contradiction. This proves the claimed bound.
\end{proof}

\section{A projective polynomial detects the nonzero values}
\begin{lemma}[Cyclotomic leading digit]\label{lem:digit}
The function $f(v)=\pi^{-t}H(v)\bmod\pi$ is nonzero and is represented on $\F_\ell^r$ by a homogeneous polynomial of degree $t$. Its nonzero set is contained in $D$.
\end{lemma}
\begin{proof}
Use coordinates $a_i,v_i\in\{0,\ldots,\ell-1\}$ in the integer exponents. Expand each $\zeta^{a_i v_i}=(1+(\zeta^{a_i}-1))^{v_i}$ in the binomial basis. Thus as a function on the coordinate grid,
\[
 H(v)=\sum_{0\le j_i<\ell} h_{\boldsymbol j}\prod_i\binom{v_i}{j_i},\quad
 h_{\boldsymbol j}=\sum_a c_a\prod_i(\zeta^{a_i}-1)^{j_i}.
\]
Every $h_{\boldsymbol j}$ is divisible by $\pi^{\sum_i j_i}$; it is also divisible by $\pi^t$, because it is an integral iterated finite difference of values of $H$ at grid points. Divide by $\pi^t$ and reduce: terms of total degree above $t$ disappear and each binomial polynomial is defined over $\F_\ell$. This represents $f$ by a polynomial of total degree at most $t$ and degree at most $\ell-1$ in each variable. It is nonzero by the definition of $t$.

For $\lambda\in\F_\ell^\times$, conjugation $\sigma_\lambda$ gives $H(\lambda v)=\sigma_\lambda(H(v))$, and $\sigma_\lambda(\pi)/\pi\equiv\lambda\pmod\pi$. Hence $f(\lambda v)=\lambda^t f(v)$. The reduced polynomial representing a function on $\F_\ell^r$ is unique, by repeated univariate interpolation. Each monomial with nonzero coefficient must therefore have degree congruent to $t$ modulo $\ell-1$. Since $f(0)=0$, none is constant. For each monomial, add $t-$degree to any positive exponent. This does not change its evaluation on $\F_\ell^r$, including when that coordinate vanishes, and makes it homogeneous of degree $t$. Finally $f(v)\ne0$ implies $H(v)\ne0$.
\end{proof}

\begin{lemma}[Projective support bound]\label{lem:prm}
A nonzero function on $\F_\ell^r$ represented by a homogeneous polynomial of positive degree $T=(\ell-1)a+b$, $1\le b\le\ell-1$, is nonzero on at least $(\ell+1-b)\ell^{r-a-2}$ projective points when $a\le r-2$, and at least one otherwise.
\end{lemma}
\begin{proof}
This is the classical minimum-distance bound for projective Reed--Muller codes \cite[Corollary 3.3]{GL}. We include an elementary proof of the lower bound to make clear that no unverified extremal-code claim is needed. First, for a nonzero reduced polynomial on $\F_\ell^n$ of degree at most $T'$, iteratively taking the highest nonzero degree in one variable shows that the number of nonzero values is at least $\prod_{i=1}^n(\ell-e_i)$ for a monomial with exponents $0\le e_i\le\ell-1$ and total degree at most $T'$. Indeed, the highest coefficient is nonzero at the inductively counted points in the remaining variables, and the resulting one-variable polynomial has at most $e_i$ roots. The smallest such product is obtained by filling exponents up to $\ell-1$ one at a time: if $0<e_i\le e_j<\ell-1$, moving one from $e_i$ to $e_j$ decreases the product. Thus for $T'=(\ell-1)a+c$, $0\le c<\ell-1$, $a<n$, the affine lower bound is $(\ell-c)\ell^{n-a-1}$; for $a\ge n$ it is one.

Now induct on $r$ for the projective statement; $r=1$ is trivial. If $a\ge r-1$, there is nothing to prove. Suppose some hyperplane misses the projective support; choose it to be $v_r=0$. Terms of the homogeneous polynomial not involving $v_r$ represent the zero function on that hyperplane and hence vanish on all vectors in their own variables. Discard those terms, divide by $v_r$, and dehomogenize at $v_r=1$. The resulting nonzero affine polynomial has degree at most $T-1=(\ell-1)a+b-1$ in $r-1$ variables. Reduction modulo $x^\ell-x$ cannot increase degree. The affine bound gives $(\ell+1-b)\ell^{r-a-2}$ nonzero points.

If every hyperplane meets the support, let $S$ be its projective size. There are $(\ell^r-1)/(\ell-1)$ hyperplanes, and each projective point lies in $(\ell^{r-1}-1)/(\ell-1)$ of them. If $a\le r-3$, apply induction to the nonzero restriction to each hyperplane to obtain
\[
 S\ge\frac{\ell^r-1}{\ell^{r-1}-1}(\ell+1-b)\ell^{r-a-3}>(\ell+1-b)\ell^{r-a-2}.
\]
If $a=r-2$, the trivial one-point lower bound per hyperplane gives $S>\ell\ge\ell+1-b$. (For $r=2$ this case uses hyperplanes that are single projective points.) This proves the claim.
\end{proof}

\begin{proof}[Proof of the lower bounds in Theorem~\ref{thm:main}]
By Lemma~\ref{lem:val}, $t\le(\ell-1)d+s$. The lower bound in Lemma~\ref{lem:prm} decreases monotonically as degree increases up to saturation at one. Lemma~\ref{lem:digit} and the fact that each projective point has $\ell-1$ nonzero vectors now give \eqref{eq:main}. In the saturated case $D\ne\varnothing$ is a union of punctured lines, giving the claimed lower bound directly.
\end{proof}

\section{Extremizers: split and nonsplit constructions}
All constructions are in the character group $X$, so they work in any realization by Dirichlet characters of the stated rank. If we construct $m\le k$ characters in each packet, we may pad both packets by identical arbitrary characters without changing $H$ or $D$.

\subsection{Split construction}
Assume $d\le r-2$ and $s\mid\ell-1$. Choose a $d$-dimensional subspace $W\subset X$, and two vectors $e_1,e_2$ independent modulo $W$. Let $S\subset\F_\ell^\times$ be the multiplicative subgroup of order $s$. Set
\[
 A=\{w+u e_1:w\in W,\ u\in S\},\qquad C=\{w+u e_2:w\in W,\ u\in S\}.
\]
These are disjoint sets of size $s\ell^d\le k$. By finite Fourier orthogonality,
\[
 H(v)=\ell^d\ind_{W^\perp}(v)\left(\sum_{u\in S}\zeta^{u x}-\sum_{u\in S}\zeta^{u y}\right),\qquad x=e_1\cdot v,\ y=e_2\cdot v.
\]
On $W^\perp$, $(x,y)$ is uniform on $\F_\ell^2$. By Lemma~\ref{lem:trace}, the bracket vanishes exactly at the origin and at the points with $x,y\ne0$ and $x/y\in S$. These are $s$ distinct projective lines in the plane, including their common origin. The complement has $(\ell-1)(\ell+1-s)$ points. Hence $\delta=\Delta_{\ell,d,s}$.

\subsection{Nonsplit construction}
Assume $s\mid\ell+1$ and $1\le s\le\ell-1$. In the two-dimensional $\F_\ell$-space $K=\F_{\ell^2}$, let $U=\{z\in K^\times:z^{\ell+1}=1\}$, the cyclic norm-one group of order $\ell+1$. Let $S\le U$ have order $s$ and take $c\in U\setminus S$. Use base packets $S,cS$ in $K$. For each $u\in cS$, the map $R_u(z)=u z^\ell$ is an $\F_\ell$-linear involution interchanging these two sets. For odd $\ell$, neither $I$ nor $-I$ is possible since Frobenius is not scalar, so its invariant dual space is a line. Each nonzero functional on that line gives identical projections of the two sets. Distinct $u,u'$ have distinct invariant dual lines: if $f=f\circ R_u=f\circ R_{u'}$, then $f=f\circ R_uR_{u'}$, but $R_uR_{u'}$ is multiplication by $u/u'\ne1$ and multiplication by $u/u'-1$ is invertible, forcing $f=0$. Thus at least $s$ projective directions agree, so the disagreement proportion is at most $(\ell-1)(\ell+1-s)/\ell^2$. The lower bound already proved, with $d=0,r=2,k=s$, shows that these are exactly the agreeing directions. For $\ell=2$ the split construction already covers the only tier and the nonsplit case is unnecessary.

Embed $K$ as a two-dimensional complement to $W$ in $X$ and replace each element of the two base packets by its whole translate of $W$. This multiplies packet size by $\ell^d$ and the disagreement density by $\ell^{-d}$, giving equality in \eqref{eq:main}.

\subsection{Saturation and examples}
If $k\ge\ell^{r-1}$, choose $W$ of dimension $r-1$ and two distinct cosets of $W$ in $X$ for the packets before padding. The Fourier difference is supported on $W^\perp\setminus\{0\}$ and nonzero there, giving $(\ell-1)/\ell^r$. For $\ell=3$ the tiers $s=1,2$ divide $2$; for $\ell=5$ the tiers $1,2,4$ divide $4$, and $3$ divides $6$. This completes the proof of Theorem~\ref{thm:main}.

\begin{center}
\begin{tabular}{@{}lll@{}}\toprule
Cubic size $k$ & Least disagreement (large rank) & Maximal agreement\\\midrule
$1$ & $2/3$ & $1/3$\\
$2$ & $4/9$ & $5/9$\\
$3$--$5$ & $2/9$ & $7/9$\\
$6$--$8$ & $4/27$ & $23/27$\\
$9$--$17$ & $2/27$ & $25/27$\\
$18$--$26$ & $4/81$ & $77/81$\\
$27$--$53$ & $2/81$ & $79/81$\\\bottomrule
\end{tabular}
\end{center}

\section{Rigidity and a rank-two reduction}
The published classification of minimum-weight projective Reed--Muller codewords \cite[Theorem 4.3]{GL} says that for homogeneous degree $(\ell-1)d+s$, $d\le r-2$, the nonzero set of a minimum word is, after a linear change of coordinates,
\begin{equation}\label{eq:geometry}
 \{v:L_1(v)=\cdots=L_d(v)=0,\quad L(v)\ne0,\quad
 M(v)\ne\omega_j L(v)\ (1\le j\le s-1)\},
\end{equation}
where $L_1,\ldots,L_d,L,M$ are independent linear forms, the $\omega_j$ are distinct elements of $\F_\ell$, and $M$ is unnecessary when $s=1$. This follows by taking the support of their displayed homogeneous minimum-word representative. We use the classification only in this section and in the extremal refinement of the analytic result.

\begin{theorem}[Equality geometry and reduced mass]\label{thm:rigidity}
Suppose $d\le r-2$ and a pair attains \eqref{eq:main}. Then $D$ has exactly the shape \eqref{eq:geometry}. After common factors are cancelled, each packet is a union, with multiplicities, of full cosets of a common $d$-dimensional subspace $W$ of $X$; the reduced size is exactly $s\ell^d$. Moreover the lower bound is attainable in this tier if and only if there exist two disjoint multisets of size $s$ in $\F_\ell^2$ whose linear projection multisets agree in exactly $s$ projective directions. Thus sharpness in all ranks reduces to an explicitly finite planar problem.
\end{theorem}
\begin{proof}
If $t<(\ell-1)d+s$, the projective lower bound at degree $t$ is strictly larger than the required value (the given degree is below saturation). Thus $t=(\ell-1)d+s$ and $\supp f=D$, since the cardinalities are equal. Apply the minimum-word classification to Lemma~\ref{lem:digit}, obtaining \eqref{eq:geometry}. Write $W^\perp=\bigcap_{i=1}^d\ker L_i$. Because $H$ is supported in $W^\perp$, Fourier inversion shows $c_{a+w}=c_a$ for all $w\in W$ (multiplying the inverse character by $\zeta^{-w\cdot v}$ changes nothing on the support). Consequently after cancellation both packets are constant on $W$-cosets, and $m$ is a multiple of $\ell^d$. If $m<s\ell^d$, applying the already-proved lower bound with $m$ instead of $k$ gives a strictly larger minimum (the tier decreases, and $d\le r-2$). Since $m\le k<(s+1)\ell^d$, we get $m=s\ell^d$.

On $W^\perp$, divide $H$ by $\ell^d$; it is the Fourier difference of the two size-$s$ packets in $X/W$. Choose a two-dimensional subspace $E$ of $W^\perp$ on which $(L,M)$ is an isomorphism onto $\F_\ell^2$ (when $s=1$, extend $L$ by any independent form). Restricting the Fourier difference to $E$ corresponds to projecting the packets to $\widehat E$. The disagreement there is exactly the complement of $s$ projective directions by \eqref{eq:geometry}. If common factors appeared on projection, the size after cancellation would be less than $s$, contradicting the lower bound for rank two. Conversely, a rank-two construction can be lifted to $X$ by the same coset replication as in Section 5. This proves both assertions.
\end{proof}

\begin{corollary}[First two tiers]\label{cor:class}
In the first tier $s=1$, every extremal pair reduces to two distinct cosets of $W$, each of size $\ell^d$. In the second tier $s=2$ for odd $\ell$, every extremal pair reduces modulo $W$ to two opposite pairs of corners of a nondegenerate parallelogram. Conversely all such pairs attain their tier bound.
\end{corollary}
\begin{proof}
The first statement follows directly from the reduced mass in Theorem~\ref{thm:rigidity}. For $s=2$, the reduced packets in $X/W$ have the forms $\{a,b\}$ and $\{c,e\}$, disjoint as multisets. Projection agrees if the two elements can be matched by one of two permutations. The agreeing set is the union of the annihilators of $\operatorname{span}(a-c,b-e)$ and $\operatorname{span}(a-e,b-c)$. Both subspaces of the dual are proper. For equality their union has relative size $2/\ell-1/\ell^2$, which is possible only for distinct hyperplanes: if either has codimension at least two, the union is at most $1/\ell+1/\ell^2<2/\ell-1/\ell^2$ for odd $\ell$; equal hyperplanes also give too small a union. Thus the two difference spans are distinct lines. Write $c=a+u,e=a+v$ with $u,v$ independent. The conditions force $b\in c+\F_\ell v$ and $b\in e+\F_\ell u$, whose intersection is $a+u+v$. The converse follows by counting the two distinct agreement hyperplanes.
\end{proof}

\begin{proposition}[Saturated equality]\label{prop:saturated}
If $d\ge r-1$, equality holds exactly when after cancellation the signed multiplicity function is constant on the cosets of some $(r-1)$-dimensional subspace $W$. On $X/W\cong\F_\ell$ it is any nonzero integral function of sum zero, with positive mass at most $k/\ell^{r-1}$.
\end{proposition}
\begin{proof}
Equality means the support of $H$ is exactly one punctured line. Its annihilator is $W$; inversion gives coset invariance as above. Conversely the Fourier transform of a nonzero integral sum-zero vector on $\F_\ell$ is nonzero at every nontrivial frequency by the cyclotomic minimal polynomial argument of Lemma~\ref{lem:trace}. This proves the characterization.
\end{proof}

For $\ell=3$ this classifies all equality pairs in every nonsaturated tier, not just their density. At $\ell=5$ the theorem gives the support and reduced mass in every tier; we do not claim a classification of all planar configurations for the $s=3,4$ tiers.

\section{Effective discrimination from the vendored Dirichlet GRH}
Fix a nonnegative smooth function $w$ supported in $(1/2,1)$ with integral one. For $g:V\to\C$ define
\[
 \widehat g(a)=\ell^{-r}\sum_{v\in V}g(v)\zeta^{-a\cdot v},\quad
 \norm{g}_A=\sum_{a\in X}|\widehat g(a)|,\quad
 \mathcal M_x(g)=\sum_{p\nmid Q}(\log p)w(p/x)g([p]).
\]
Here $[p]$ is the evaluation vector of the prime. The smooth character estimate in the article (\texttt{rec:lem:gram}) gives, using Dirichlet GRH,
\begin{equation}\label{eq:smooth}
 \big|\mathcal M_x(g)-x\E_Vg\big|\le C_w\norm{g}_A\sqrt x\log(2Q)\qquad(x\ge 2),
\end{equation}
with $C_w$ independent of $\ell,r,Q,g$. For completeness, apply Mellin inversion to $\sum_n\Lambda(n)\eta(n)w(n/x)$ for the primitive inducer $\eta$ of a character modulo $Q$. The pole of $\zeta$ contributes $x$ for the principal character. The residues at nontrivial zeros contribute $-\sum_\rho x^\rho\widetilde w(\rho)$, where $\widetilde w(s)=\int_0^\infty w(y)y^{s-1}\,dy$ decays rapidly on vertical lines. On GRH, the real parts are $1/2$, so the standard unit-interval zero count $O(\log(2f(2+|t|)))$ for conductor $f\le Q$ bounds this sum by $O_w(\sqrt x\log(2Q))$. Shifting the contour to $-1/4$ and using the functional equation bounds the remaining integral by $O_w(x^{-1/4}\log(2Q))$, with a bounded contribution at $s=0$. Higher prime powers cost $O_w(\sqrt x)$ by the Chebyshev bound and primes dividing $Q$ cost $O_w(\log Q)$. Summing the resulting character estimates against the Fourier coefficients gives \eqref{eq:smooth}. This is precisely where the source GRH input is used. See also \cite{Davenport,Bellaiche} for the classical analytic and Fourier framework.

\begin{theorem}[Rank-independent quantitative detection]\label{thm:effective}
For every distinct pair of size $k$, with reduced size $m\le k$,
\begin{equation}\label{eq:eff}
 \mathcal M_x(\ind_D)\ge \frac{\ell-1}{\ell m}x-2C_w\frac{\ell-1}{\ell}\sqrt x\log(2Q).
\end{equation}
Under Dirichlet GRH, in particular, some prime $p\in(x/2,x)$ with $p\nmid Q$ distinguishes the two Euler factors whenever $x\ge K_w k^2\log^2(2Q)$, with an absolute constant for fixed $w$. Taking $K_w$ sufficiently large, the number of distinguishing primes in that interval is $\gg_w x/(k\log x)$, uniformly in $\ell,r,Q$.
\end{theorem}
\begin{proof}
Define a real nonnegative test
\[
 g(v)=\frac1{2\ell m^2}\sum_{j=1}^{\ell-1}|H(jv)|^2.
\]
By cyclic Parseval and $\sum b_i=0$, the numerator sum is $\ell\sum_i b_i(v)^2$. The positive and negative masses of $b_i(v)$ are at most $m$, so $\sum b_i^2\le2m^2$. Thus $0\le g\le\ind_D$. Finite Fourier Parseval gives
\[
 \E_Vg=\frac{\ell-1}{2\ell m^2}\sum_a c_a^2\ge\frac{\ell-1}{\ell m},
\]
since $c_a$ are integers. Also the Fourier norm is submultiplicative, invariant under conjugation and $v\mapsto jv$, and $\norm H_A=2m$, so $\norm g_A\le2(\ell-1)/\ell$. Apply \eqref{eq:smooth}. If $\sqrt x\ge4C_w k\log(2Q)$, the right-hand side of \eqref{eq:eff} is at least $(\ell-1)x/(2\ell k)$; then divide by $\|w\|_\infty\log x$ to obtain the count. Constants can be enlarged to ensure $x\ge2$ and strict positivity. The same argument without prime sums proves the unconditional universal estimate $\delta\ge(\ell-1)/(\ell m)$, weaker than the staircase except at its first-tier lower endpoint and at some special cases.
\end{proof}

\begin{corollary}[Full density recovery for extremal pairs]\label{cor:measure}
For a pair attaining the bound \eqref{eq:main},
\[
 \mathcal M_x(\ind_D)=x\delta+O_w(2s\sqrt x\log(2Q)).
\]
For a saturated equality case, the same holds with $2$ in place of $2s$. In particular for $\ell=3,5$, extremal density is measured with a fixed relative error at $x\gg_w k^2\log^2(2Q)$, with no dependence on the rank. For $d\le r-2$, if it is known that $\delta$ is one of the $\ell-1$ tier values $\Delta_{\ell,d,j}$ and the pair has the corresponding equality-support geometry, the exact value can be recovered by rounding at $x\gg_{w,\ell} \ell^{2d}\log^2(2Q)$.
\end{corollary}
\begin{proof}
By Theorem~\ref{thm:rigidity}, in $W^\perp$ the disagreement set is the complement of $s$ lines through the origin in a two-dimensional quotient (for $s=1$, one may use any independent extra coordinate). On $\F_\ell^2$ its indicator is $1-\sum_{i=1}^s\ind_{L_i}+(s-1)\ind_{\{0\}}$. Each subgroup indicator has Fourier norm one; pulling it back by a surjective linear map preserves that norm, and multiplication by $\ind_{W^\perp}$ does not increase it. Hence $\norm{\ind_D}_A\le2s$. In the saturated case it is a difference of two subgroup indicators. Apply \eqref{eq:smooth}. For fixed $\ell$, $\delta\gg_\ell1/k$ and adjacent tier values differ by $(\ell-1)/\ell^{d+2}$; the asserted measurement and rounding estimates follow. The rounding statement explicitly assumes the promised geometry or corresponding norm bound; it is not a claim about arbitrary pairs or arbitrary near-minimum densities.
\end{proof}

\section{Consequences, checks, and further questions}
The construction supplies pairs of distinct Euler products that are locally indistinguishable at nearly all primes, while the lower bound proves they cannot agree more often than the stated threshold. For cubic and quintic products the maximal agreement is determined exactly by size and rank. This is a finer invariant than the simple uncertainty estimate $\abs{\supp H}\abs{\supp c}\ge\ell^r$: the latter would only yield $\delta\ge1/(2k)$. The staircase uses integer multiplicity and cyclotomic arithmetic essentially. It also identifies a structural reason why agreement plateaus persist over large ranges of packet sizes.

The supplied \texttt{verify.py} independently computes disagreement by exact integer projection counts (not by floating-point roots of unity), enumerates small packet cases, checks the split and nonsplit constructions, and verifies the valuation bound for many small signed count vectors. Its output is supporting evidence, not a substitute for the proof. The source hashes and the inspected files are recorded in \texttt{sources.md}. A \texttt{README.md} describes the scope and reproduction steps.

\paragraph{Further research.} The first tier not settled by our divisibility criterion is $\ell=7,s=5$; Theorem~\ref{thm:rigidity} reduces sharpness to a planar projection problem. For larger primes, classify the $s$ admitting two $s$-point multisets with $s$ agreeing directions; if the lower bound is not attained, find the next support size. A second problem is to determine the entire high-agreement spectrum above the minimum, analogous to the article's Kasami--Tokura analysis at exponent two. Third, seek a stronger arithmetic input that reduces the $k^2\log^2 Q$ detection scale or yields matching lower bounds for worst-case moduli. Fourth, extend the augmentation and projective-code method to character groups of prime-power exponent, where local polynomials are no longer determined by one trace, or to unequal packet sizes. Finally, formalizing the augmentation lemma, Newton bound, polynomial support argument and the analytic estimate in the pinned Lean package would make the dependency on the five exported source theorems fully machine-checkable.

\begin{thebibliography}{9}
\bibitem{vendor} ProveIt, \emph{Vendored grh-lean (readable edition)}, \url{https://github.com/VladimirReshetnikov/ProveIt/tree/main/Analysis/GeneralizedRiemannHypothesis/vendor/grh-lean}. Upstream pinned commit \texttt{72fead102f663d5c7a407}\allowbreak\texttt{964a13dc0cf9d977207}. See \texttt{grh/README.md}, \texttt{src/Statements.lean}, \texttt{src/Main.lean}, \texttt{src/New/Main.lean}.
\bibitem{article} ProveIt, \emph{Beyond the Riemann hypothesis for Dirichlet L-functions: consequences, transfers and limits}, merged article, \url{https://github.com/VladimirReshetnikov/ProveIt/tree/main/Analysis/GeneralizedRiemannHypothesis/Article}. In particular \texttt{sections/rec.tex}, labels \texttt{rec:lem:gram}, \texttt{rec:thm:qmatch}, \texttt{rec:thm:qmatchfinite}, \texttt{rec:sec:open}. Blob \texttt{da56d6ad5a1a699e335dcfd327bcdc9492c64017}.
\bibitem{GL} S. R. Ghorpade and R. Ludhani, On the minimum distance, minimum weight codewords, and the dimension of projective Reed--Muller codes, \emph{Adv. Math. Commun.} \textbf{18} (2024), 360--382. \href{https://doi.org/10.3934/amc.2023035}{doi:10.3934/amc.2023035}. Corollary 3.3 and Theorem 4.3.
\bibitem{Sorensen} A. B. S\o rensen, Projective Reed--Muller codes, \emph{IEEE Trans. Inform. Theory} \textbf{37} (1991), 1567--1576. See also the author's correction, \href{https://arxiv.org/abs/2310.03574}{arXiv:2310.03574} (2023).
\bibitem{Davenport} H. Davenport, \emph{Multiplicative Number Theory}, 3rd ed., revised by H. L. Montgomery, Springer, 2000.
\bibitem{Bellaiche} J. Bella\"iche, \emph{Th\'eor\`eme de Chebotarev et complexit\'e de Littlewood}, 2016 preprint (Fourier complexity and effective Chebotarev).
\end{thebibliography}
\end{document}
